\section{Isolation of the tripled curve}\label{sec:neighborhood}

Let $D$ have genus $h=3$ or $4$, and let $N$ be a general stable
bundle of rank two with determinant $K_D$. Write
$Y=\operatorname{Tot}_D N$. By Theorem~\ref{thm:classification},
the degree-three cycles in $Y$ away from $3[D]$ are smooth embedded
curves in finitely many orbits of fiber dilation. Each orbit comes
from a line subbundle $S\subset N$ of degree one. For a generator
$s\in H^0(D,S^3)$ its members are the curves
\[
                 F_a=\{y^3=a s\}\subset\operatorname{Tot}_D S,
                      \qquad a\in\mathbb C^\times.
\]
The three zeros of $s$ are simple, and $H^0(S)=H^0(S^2)=0$.
We first compute the obstruction to deforming $F_a$ when the
threefold varies.

\subsection{The normal bundle of a cubic}

It suffices to take $a=1$. Put $F=F_1$, write $f:F\to D$ for the
projection and $Q=K_DS^{-1}$, and denote the class of the extension
of $Q$ by $S$ defining $N$ by $e\in H^1(SQ^{-1})$.
The equation of $F$ in $\operatorname{Tot}_D S$ gives the normal
sequence
\begin{equation}\label{eq:cubic-normal-sequence}
 0\longrightarrow f^*S^3\longrightarrow N_{F/Y}
               \longrightarrow f^*Q\longrightarrow0 .
\end{equation}
Its connecting map on global sections is cup product with
$3y^2f^*e$, up to sign. Indeed, local splittings of
$N\to Q$ lift a $Q$-displacement to vertical vector fields whose
differences are the corresponding $S$-displacements given by $e$.
The derivative of the cubic equation on these differences is
multiplication by $3y^2$.

\begin{lemma}\label{lem:cubic-normal-cohomology}
For all the cubics of Theorem~\ref{thm:classification}, after shrinking
the dense open of normal bundles if necessary, the tangent
$H^0(F,N_{F/Y})$ is generated by varying $a$, and projection in
\eqref{eq:cubic-normal-sequence} induces an isomorphism
\[
        H^1(F,N_{F/Y})\ \xrightarrow{\ \sim\ }\ H^1(F,f^*Q)
           = H^1(D,QS^{-2}) \cong H^0(D,S^3)^\vee .
\]
\end{lemma}

\begin{proof}
Use $f_*\mathcal O_F=\bigoplus_{i=0}^2 S^{-i}$, with the $i$th
summand represented by $y^i$ in a local frame. The three spaces
$H^0(QS^{-i})$ have dimensions $h-2,h-3,h-3$ by Riemann--Roch.
The first maps isomorphically to $H^1(S)$ in the $y^2$ summand of
$H^1(f^*S^3)$. This is three times the connecting map for $N$,
which is acyclic by degree-two rigidity
\cite[Theorem~2.2]{BP06}. In genus three these are the only nonzero
blocks.

In genus four take generators $q_1$ and $q_2$ of $H^0(QS^{-1})$
and $H^0(QS^{-2})$. Their blocks map to $H^1(S^3)$ and $H^1(S^2)$,
respectively. Since $y^3=s$, both Serre pairings are, up to the
common factor and sign, $\langle e,sq_1q_2\rangle$. The product is
a nonzero section of $Q^2$, the Serre dual of the extension space.
Its annihilator is a hyperplane in each projective space
$\mathbb P H^1(SQ^{-1})$ of dimension six. The possible $S$ with
a cubic section range over a three-dimensional locus in
$\operatorname{Pic}^1D$. These exceptional extensions therefore
have image of dimension at most eight in the nine-dimensional
moduli space of bundles with fixed determinant.

The connecting map in \eqref{eq:cubic-normal-sequence} is consequently
an isomorphism in both genera. The remaining tangent space is
$H^0(f^*S^3)=H^0(S^3)$, spanned by the derivative of $F_a$.
Serre duality gives the asserted obstruction space, since its other
two summands are dual to $H^0(S)$ and $H^0(S^2)$.
\end{proof}

\subsection{The second infinitesimal neighborhood}

Let $D\subset X$ be an embedding in a smooth quasi-projective
threefold, with an identification $N_{D/X}=N$. Write
$\mathcal X\to\mathbb A^1$ for its deformation to the normal cone, obtained by
blowing up $D\times0$ in $X\times\mathbb A^1$ and removing the proper
transform of $X\times0$. The central fiber is $Y$ and every
noncentral fiber identifies with $X$. For an embedded cubic
$F\subset Y$, let
\begin{equation}\label{eq:cubic-unfolding}
        \kappa_F(X)\in H^1(F,N_{F/Y})
\end{equation}
be the restriction of the Kodaira--Spencer class of this deformation
at zero, projected to the normal bundle of $F$. This is the
obstruction to lifting $F$ over the tangent vector $\partial_t$ of
the parameter line.

The class depends only on the second infinitesimal neighborhood
defined by $\mathcal I_D^3$. Indeed, in adapted coordinates $(z,n)$
along $D$, write a change of charts as
\[
  z'=h_0(z)+A(z)n+O(n^2),\qquad
  n'=G(z)n+B(z)(n,n)+O(n^3).
\]
In normal-cone coordinates replace $n$ by $tn$ and divide $n'$ by
$t$. The derivative at zero has horizontal component $A(z)n$
and vertical component $B(z)(n,n)$. Projecting this cocycle to
$N_{F/Y}$ gives \eqref{eq:cubic-unfolding}. The definition thus
applies to embeddings without a retraction onto $D$. Dilation
lifts to the normal-cone family, acting on its parameter by the
inverse character. Nonvanishing on one member of a dilation orbit
is equivalent to nonvanishing on every member.

\begin{proposition}\label{prop:isolation}
For $h=3,4$, suppose that $\kappa_F(X)\ne0$ for one cubic in
each orbit of Theorem~\ref{thm:classification}. Then $3[D]$ is
isolated among all compact effective cycles of local degree three
in a tube about $D$ in $X$. This condition on the second
infinitesimal neighborhood is Zariski open in algebraic families
with identified $D$ and $N$. It is attained by smooth
quasi-projective threefolds with normal bundle $N$ and
determinant-preserving changes of the fiber coordinates.
\end{proposition}

\begin{proof}
In a family with fixed $D$ and $N$, the restricted
Kodaira--Spencer classes in \eqref{eq:cubic-unfolding} vary
algebraically in the fixed obstruction lines. There are only
finitely many conditions. We exhibit an algebraic family where
each vanishing condition cuts out a proper hyperplane.

Start with $Y$. Fix a point $p\in D$, a frame of $N$ on an affine
neighborhood $U$, and a local parameter $u$ regular on $U$ with
no other zero there. Choose a covector $\ell$ on the framed plane
which is nonzero on every one of the finitely many lines $S_p$,
and a nonzero $v_\ell\in\ker\ell$. Reglue the plane bundle over
$U$ to its restriction over $D\setminus\{p\}$ by composing the
transition on the overlap with
\begin{equation}\label{eq:normal-quadratic-shear}
        n\longmapsto n+v_\ell\,\ell(n)^2
                         \sum_{j=1}^m c_j u^{-j}.
\end{equation}
This polynomial automorphism has determinant one and its inverse
negates the coefficients. It fixes the zero section to first order.
The resulting smooth threefold is affine of finite type over $D$,
hence quasi-projective. Its projection to $D$ has the same vertical
tangent determinant as the original bundle, since the added
transition has determinant one.

For a given $S$ the restriction of $v_\ell\ell^2$ to $S^2$,
followed by projection to $Q$, generates $QS^{-2}$ at $p$.
The Kodaira--Spencer class on the cubic therefore projects to the
principal-parts boundary in $H^1(QS^{-2})$, in the $y^2$ summand of
$H^1(f^*Q)$. For sufficiently large $m$, these boundaries span
that space by $H^1(QS^{-2}(mp))=0$. Multiplying the principal
parts by the local generator preserves their span. Choose $m$
large enough for all the $S$. By Lemma~\ref{lem:cubic-normal-cohomology}
each $\kappa_F$ is then a nonzero linear functional of the $c_j$,
with values in a one-dimensional space. Arbitrarily small choices
outside the union of their kernels satisfy the hypothesis.

Now fix an $X$ satisfying this hypothesis. At a nonzero cubic $F$
the tangent sequence for the relative Hilbert scheme of
$\mathcal X/\mathbb A^1$ is exact through
\[
 0\longrightarrow H^0(N_{F/Y})\longrightarrow T_{[F]}\operatorname{Hilb}
   \longrightarrow T_0\mathbb A^1
   \xrightarrow{\ \kappa_F\ }H^1(N_{F/Y}).
\]
Indeed, local lifts of the embedding over a first-order parameter
differ on overlaps by precisely the normal cocycle
\eqref{eq:cubic-unfolding}. The Hilbert tangent has dimension one.
The actual orbit $F_a$ has dimension one in the central fiber
and its derivative generates that tangent. Equality of the local
dimension and embedding dimension gives a regular local ring.
Thus the relative Hilbert scheme near $[F]$ consists entirely of
the central orbit. In particular this conclusion concerns every
deformation of the embedded curve.

Suppose extra cycles converged to $3[D]$ in $X$. In a smooth
tubular identification let $t_i>0$ be the ratio of the maximum
normal distance of the $i$th cycle to a fixed small radius. Then
$t_i\to0$. Regard the cycles in the fibers $\mathcal X_{t_i}$,
using the divided normal coordinates. They lie in a common
closed disk and reach its prescribed radius. The smooth structures
after this rescaling converge on a larger disk to the complex
structure of $Y$. A closed K\"ahler form on the central disk tames
these structures for large $i$. Its integral on every cycle is
three times its integral on $D$ by retraction to the zero section.
It follows that the masses are uniformly bounded. Cycle compactness,
including Hausdorff convergence of support, gives a limit of degree
three meeting the prescribed positive radius
\cite[Theorem~2.11 and Remark~2.13]{DIW}. Theorem~\ref{thm:classification}
identifies the limit with a single smooth cubic of multiplicity one.

Openness of the embedded locus and its topology in the cycle space
give smooth embedded nearby curves of the same genus
\cite[Proposition~2.21]{DIW}. To pass to the relative Hilbert
scheme, use a projective completion of the normal-cone family.
The ample degree of the cycles is locally constant, so the Hilbert
polynomial of the nearby embedded curves agrees with that of $F$.
Any flat Hilbert limit has support $F$ and generic multiplicity
one by convergence of supports and cycles. The kernel of its
structure sheaf onto $\mathcal O_F$ has finite length, hence
vanishes by equality of Hilbert polynomials. The resulting Hilbert
limit $[F]$ contradicts the central-orbit description above.
\end{proof}
